\section{Results}
\label{sec:taxonomies}

The reviewing process resulted in a dataset containing responses to 21 question items on 445 benchmark articles, annotated by 29 experts in the areas of NLP and machine learning. \Cref{fig:results_overview} shows that the number of included articles increases significantly in each subsequent year. The dataset contains information covering all of the stages of the benchmarks, from how they initially define their phenomenon of interest, to which tasks they select in an attempt to measure this phenomenon, to the metrics they use to estimate and compare the performance of language models on these tasks, to the claims they make about their benchmark's ability to accurately measure the phenomenon. \Cref{fig:results_sankey} shows the results of key questions from the reviewing codebook which motivate our recommendations.

\begin{figure}[t!]
    \centering
    \includegraphics[width=\linewidth]{images/benchmark_review_results_overview.pdf}
    \caption{\textbf{Summary of reviewed articles.} (A) Three most common categories of benchmark phenomena, grouped into general capabilities, general applications, and specific applications. (B) Number of articles by publication year and number which discuss the construct validity of their benchmark.}
    \label{fig:results_overview}
\end{figure}



\paragraph{\textcolor{phenomenon}{Phenomenon}} 
The reviewed benchmarks cover a wide range of phenomena (\Cref{fig:results_overview} (A)), including areas such as reasoning (18.5\%), alignment (8.1\%) and code generation (5.7\%). 
% taxolomy root:
% NLP                      20.000000
% Reasoning                18.461538
% Agents                    8.791209
% Alignment                 8.131868
Most of the articles provided a definition of their measured phenomenon (78.2\%). Of the articles that provided definitions, 52.2\% of these definitions are widely agreed upon, but 47.8\% are contested, addressing phenomena with many possible definitions or no clear definition at all (\Cref{fig:results_sankey}).
For example, a benchmark measuring the extent to which LLM generations correspond to established psychometric categories  \cite{renValueBenchComprehensivelyEvaluating2024} has clear, widely agreed definitions. However, as a category, alignment benchmarks often target phenomena with contested definitions (e.g. `harmlessness').

The definitions of the phenomena also varied in whether they defined the phenomenon they tested as a composite (61.2\%) or a single unified whole  (36.5\%). For example, some phenomena can be tested alone, (e.g. measuring the ability to traverse a 2D map~\cite{nasirGameTraversalBenchmarkEvaluatingPlanning2024a}), while other phenomenon are overarching abilities integrating many sub-abilities (e.g., a model's `agentic capabilities' requiring sub-abilities such as intent recognition, alignment, and structured output generation~\cite{liuAgentBenchEvaluatingLLMs2024}).

% definition_integrity
% Composite phenomenon               61.233480
% Single cohesive phenomenon         36.563877
% Authors' description is unclear     2.202643

% what is the most common task -> task_ecology_clean
% percentage of tasks from benchmarks and/or exams task_source_clean
% Constructed	40.659341
% Representative	36.923077
% Partial	32.307692
% Complete	9.230769

% Constructed	40.659341
% Representative	36.923077
% Partial	32.307692
% Complete	9.230769
\paragraph{\textcolor{task}{Task}} 
The tasks chosen to measure the target phenomena varied widely, ranging from answering medical licensing exam questions~\cite{ouyangCliMedBenchLargescaleChinese2024a} and detecting errors in computer code~\cite{shahStackEvalBenchmarkingLlms2024a} to reconciling conflicting information on Wikipedia~\cite{houWikiContradictBenchmarkEvaluating2024a}. Less than 10\% of benchmarks used complete real-world tasks, such as writing a correct SQL query given a natural language query and a database structure~\cite{changDrspiderDiagnosticEvaluation2023}.
Overall, 40.7\% of all reviewed benchmarks make use of constructed tasks, such as reading fictional multi-party conversations and answering questions about the beliefs of the conversation participants to test `theory of mind'~\cite{kimFANToMBenchmarkStresstesting2023}, with 28.5\% using exclusively constructed tasks. 
Partially real-world tasks, such as accomplishing e-commerce tasks collected from real people on a mock website \cite{yaoWebShopScalableRealworld2022a}, and representative tasks, such as answering exam-style science questions \cite{wangSciBenchEvaluatingCollegelevel2024}, are used in 32.3\% and 36.9\% of reviewed benchmarks, respectively.

% task used at least partially:
% Constructed	40.659341
% Representative	36.923077
% Partial	32.307692
% Complete	9.230769

% 	task used exclusivley:	
% Constructed	28.571429	
% Representative	23.076923	
% Partial	21.758242	
% Complete	7.912088	

\begin{figure}
    \centering
    \includegraphics[width=\linewidth]{images/results_sankey.pdf}
    \caption{\textbf{Key codebook results.} The distribution of codebook responses on selected items. In each column, the options are ordered from most to least preferred for high construct validity. The shaded area indicates the benchmarks that follow the best practices for all five items.}
    \label{fig:results_sankey}
\end{figure}

Benchmarks included task items from various sources, with only 33.6\% relying on a single source. Authors most commonly handcrafted new task items (43.3\%), followed by reusing data from existing benchmarks (42.6\%) and generating data with LLMs (31.2\%). Human exams and other pre-existing sources were used in 38.2\% of benchmarks. Among those with a single task source, the most common was another benchmark (7.7\% of benchmarks sourced their tasks solely from another benchmark).

% using as part of process:
% Author-crafted	43.296703
% Another benchmark	42.637363
% LLM-generated	31.208791
% Real task	29.670330
% Procedurally-generated	26.813187
% Crowd-sourced	19.120879
% Expert-crafted	12.967033
% Human exams	9.450549
% Unknown	0.879121

% only using:
% Another benchmark         7.692308
% Real task                 6.813187
% Author-crafted            6.593407
% Procedurally-generated    3.076923
% Expert-crafted            2.857143
% LLM-generated             2.417582



% Criterion sampling: creators take items from a larger set based on specified rules
% example: filtering recent ML papers to produce literature review questions (ajithLitSearchRetrievalBenchmark2024)
% Targeted sampling: creators define a task space and choose tasks within it strategically
% example: compiling news sources with differing opinions and asking the LLM to synthesise the information (huangEmbraceDivergenceRicher2024)
% Random sampling: creators define a task space and randomly sample from it
% example: collect samples of informal text from the Semantic Scholar Corpus and evaluate rewriting in academic prose (diaoDoolittleBenchmarksCorpora2023)

Task items within any benchmark are effectively a sample from a much larger conceptual set of possible items that could be used to operationalise the phenomenon. The methodology used to select this sample significantly impacts the benchmark's validity. In 12.3\% of cases, authors exclusively used readily accessible datasets as the source of their task items, a practice known as \emph{convenience} sampling \cite{alvi2016manual}. Another 27.0\% incorporated convenience sampling as part of their sampling strategy. Overall, 55.2\% of benchmarks used at least partially \emph{targeted} sampling, which involved defining a task space and strategically selecting items within it (e.g., compiling news sources with differing opinions and testing how well models can synthesise the information~\cite{huangEmbraceDivergenceRicher2024}). 46.2\% used \emph{criterion} sampling, where items are selected from a larger set based on specified rules, as part of their strategy (e.g., filtering recent machine learning articles to produce literature review questions~\cite{ajithLitSearchRetrievalBenchmark2024}). Finally, 17.1\% used \emph{random} sampling, which involves defining a task space and randomly selecting items from it (e.g., collecting samples of informal text from Semantic Scholar and evaluate rewriting in academic prose~\cite{diaoDoolittleBenchmarksCorpora2023a}). Of all reviewed articles, 17.5\% exclusively used targeted sampling, 10.8\% exclusively used criterion sampling, and 7.0\% exclusively used random sampling. 

% TO ADD: GIVE MORE INFO ON OTHER SAMPLING TECHNIQUES AND WHAT THEY MEAN, BASED ON THE DATA BELOW:

% partially using sampling method:
% Targeted	55.164835
% Criterion	46.153846
% Convenience	39.340659
% Random	17.142857
% Unknown	4.395604

% only using sampling method:
% Targeted       17.582418
% Convenience    12.307692
% Criterion      10.769231
% Random          7.032967

%{\color{red} TO DO: SIZE OF BENCHMARKS INFO}

% 4 most common response formats: response_format_clean: “structured response” is useful for the auxiliary tasks
	
% Free response	46.373626
% Multiple choice	40.000000
% Short free response	38.461538
% Structured	21.098901
% Interaction	6.593407
% Logits	1.318681

% what metric was featured how much: metric_definition_clean
% Exact match	81.318681
% Soft match	20.879121
% LLM-as-a-Judge	17.142857
% Human ratings	12.967033
% Distribution	12.087912
% LLM post-processing	9.450549
% Reward	8.791209
% Correlation	5.054945


Benchmarks included task items that require responses in a range of formats. Free response was the most common response format (used partially by 46.4\%, exclusively by 17.8\%), followed by multiple choice (used partially by 40.0\%, exclusively by 18.5\%). Short free response, where responses were limited to a few words, was used partially by 38.5\% and exclusively by 13.2\%, and other structured response formats such as JSON  were used partially by 21.1\% and exclusively by 8.6\%.

% response format used at least aprtially:
% Free response	46.373626
% Multiple choice	40.000000
% Short free response	38.461538
% Structured	21.098901
% Interaction	6.593407
% Logits	1.318681

% response format used exclusivley:
% Multiple choice        18.461538
% Free response          17.802198
% Short free response    13.186813
% Structured              8.571429
% Interaction             2.197802
% Logits                  1.098901

\paragraph{\textcolor{metric}{Metric}} 
The most common metric used to score the benchmarking tasks was exact matching (used at least partially by 81.3\%, exclusively by 40.7\%). Other commonly used metrics include soft match scores, which have an exact correct answer but allow for partial credit (used at least partially by 20.9\%, exclusively by 0.9\%), LLM-as-a-judge (at least partially by 17.1\%, exclusively by 3.1\%), and human ratings (at least partially by 13.0\%, exclusively by 1.8\%). Once the responses were scored, 16.0\% used uncertainty estimates or statistical tests to compare the results.

% percentage at least partially using metric:
% Exact match	81.318681
% Soft match	20.879121
% LLM-as-a-Judge	17.142857
% Human ratings	12.967033
% Distribution	12.087912
% LLM post-processing	9.450549
% Reward	8.791209
% Correlation	5.054945

% percentage only using metric
% Exact match            40.659341
% Reward                  3.296703
% LLM-as-a-Judge          3.076923
% Distribution            2.857143
% Human ratings           1.758242
% Soft match              0.879121
% Correlation             0.439560
% LLM post-processing     0.219780


% percent that compare results: results_comparison
% No     64.757709
% Yes    35.242291
\paragraph{\textcolor{claims}{Claims}}

To support their results, 53.4\% of articles presented evidence for the construct validity of their benchmark. This included benchmarks using real-world tasks to ensure the problems reflected actual coding~\cite{huangDAcodeAgentData2024} or authors suggesting that their designs better captured user intents~\cite{wangUsercentricMultiintentBenchmark2024}. 35.2\% of articles also included a comparison to other benchmarks of similar phenomena, 32.4\% to a human baseline, 31.2\% to a more realistic setting, enabling an understanding of similarities and differences.